\chapter{工程化实现与管理建议}

\section{从模型到产线系统的映射}
理论优化结果需要映射为可执行规则。建议将最优策略参数化为三类配置：
\begin{enumerate}
  \item \textbf{基础配置}：检查节点、换刀阈值、复检开关；
  \item \textbf{质量配置}：抽检标准、判定容差、异常处理流程；
  \item \textbf{数据配置}：寿命数据更新频率、统计窗口长度。
\end{enumerate}

\section{在线更新机制}
考虑刀具批次变化和设备老化，建议采用滚动更新：
\[
\Theta_{t+1}=\mathcal U(\Theta_t,\mathcal D_t),
\]
其中 $\mathcal D_t$ 为最近窗口的生产数据。每次更新后重新求解策略，并进行“新旧策略切换收益评估”：
\[
G_t=\phi(\pi_t)-\phi(\pi_{t+1}).
\]
仅当 $G_t$ 超过阈值才执行策略替换，避免频繁调参。

\section{异常工况处理}
若出现连续异常（如短时间内多次故障），建议触发保护机制：
\begin{itemize}
  \item 将检查间隔临时缩短到安全模式；
  \item 强制执行提前换刀；
  \item 启动设备状态诊断并记录异常标签。
\end{itemize}
该机制可显著降低极端场景下的质量损失。

\section{质量-成本协同看板}
建议建立以策略为中心的可视化看板，核心指标包括：
\begin{enumerate}
  \item 当日单位合格品成本；
  \item 废品率与误停率；
  \item 检查执行率与策略偏离度；
  \item 预测收益与实际收益差异。
\end{enumerate}
该看板既可支撑班组管理，也可用于策略回溯评估。

\section{实施路线建议}
按照“试点 -- 验证 -- 推广”三阶段实施：
\begin{enumerate}
  \item \textbf{试点阶段}：单设备上线策略，验证数据闭环；
  \item \textbf{验证阶段}：跨班组对比，评估鲁棒性；
  \item \textbf{推广阶段}：标准化参数模板，形成 SOP。
\end{enumerate}

\section{本章小结}
本章将优化模型转化为可执行流程，给出在线更新、异常处理、看板监控和分阶段部署建议，使作业结论更接近真实工业场景。
